\section{Прогнозні інтервали, детекція та задача визначення розміру}
\label{sec:sizing}

Робочий процес поєднує імовірнісне обмеження на пропускну здатність із
необов'язковим обмеженням на детекцію; підключається будь-який детектор, що дає
виміряні $(\pd,\pf,\rho)$.

\subsection{Означення та припущення}

\begin{Def}[Клас відмов і детектор]
\label{def:faults}
Підмножина $B$ розміру $\kf$ є \emph{несправною}, якщо кожен її член незалежно
пропускає раунд пропозиції з імовірністю~$\pd$, тоді як справні валідатори
пропускають його з імовірністю $\pf<\pd$ (клас omission; equivocation
виключено). Детектор позначає валідатор~$i$, коли його частка пропусків
протягом $\w$ раундів перевищує $\thr\in(\pf,\pd)$; детекція успішна, якщо
позначено хоча б одного члена~$B$.
\end{Def}

\begin{Assumption}[Обмінна кореляція]
\label{as:rho}
Індикатори участі в~$B$ обмінні з попарною кореляцією $\rho\in[0,1)$, тож
середнє несе design effect $1+(\kf-1)\rho$ у сенсі Kish'а~\cite{kish1965}.
\end{Assumption}

\subsection{Прогнозні інтервали та межа детекції}
\label{sec:detection}

З регресором $\mathbf{z}=(1,\log\cc,-\log\n)^{\!\top}$ і
$\hat\vartheta=(\log\Tzero,\gamma,\beta)^{\!\top}$ метод
дельти~\cite{seber2003} дає прогнозну дисперсію
$\hat{V}=\mathbf{z}^{\!\top}\widehat{\Cov}(\hat\vartheta)\mathbf{z}+\hat\sigma^{2}$.
Наближений $(1-\epsp)$-прогнозний інтервал для $\TPS(\n,\cc)$ тоді
$\exp\{\log\widehat{\TPS}\mp\zq_{1-\epsp/2}\hat{V}^{1/2}\}$;\label{st:pi}
односторонню нижню межу $\underline{\TPS}_{1-\epsp}$ беруть із $\zq_{1-\epsp}$.
Залишковий бутстреп~\cite{efron1993} дає емпіричні квантилі, коли нормальність
залишків сумнівна. Покриття перевіряють на відкладених значеннях~$n$ у межах
каліброваного діапазону (розділ~\ref{sec:exp}).

При $\kf=\lceil\phi\n\rceil$ евристика design effect дає наближену ймовірність
пропуску\label{st:rho}
\begin{equation}
\pmiss(\n,\w)\lesssim\exp\!\Bigl\{-\frac{2\w\kf(\pd-\thr)^{2}}{1+(\kf-1)\rho}\Bigr\},
\label{eq:pmiss}
\end{equation}
яка незростає за $\n$ і $\w$. Поріг $\pmiss\le\epss$ дає
\begin{equation}
\nmin=\left\lceil
\frac{1+(\kf-1)\rho}{\phi}\cdot
\frac{\log(1/\epss)}{2\w(\pd-\thr)^{2}}\right\rceil .
\label{eq:nmin}
\end{equation}
Вираз~\eqref{eq:pmiss} --- інженерне наближення, що адаптує
Hoeffding~\cite{hoeffding1963} через Kish~\cite{kish1965}
($\neff=\kf\w/(1+(\kf-1)\rho)$), а не нова точна межа концентрації.

\paragraph{Виміряний детектор.}
X5 дає частоти позначення окремих валідаторів за оптимального за Юденом порогу
участі~$0.90$: $\hat p_{d}=\resPd$ для несправних і $\hat p_{f}=\resPf$ для
справних валідаторів ($J=\resYouden$). Підстановка цих значень
у~\eqref{eq:nmin} з $\thr$ посередині між ними, $\rho=0$, $\phi=0.2$, $\w=30$ і
$\epss=0.01$ дає $\nmin\approx\resNminDet$, що значно більше за будь-яке
допустиме за пропускною здатністю~$\n$. Отже, з цим детектором детекційне
обмеження невиконуване, тож формально допустима множина порожня; ми його
послаблюємо (не застосовуємо): звітовані вікна використовують
$\nmin=4$ ($3f+1$ при $f=1$) і визначаються лише пропускною здатністю.
Програма з двома обмеженнями застосовна до детекторів із $J$, відділеним від
нуля; її компроміс тут не продемонстровано.

\subsection{Програма визначення розміру}

\begin{Problem}[Визначення розміру множини валідаторів]
\label{prob:sizing}
За пікового попиту $\Dpk$, рівнів $\epsp,\epss$, вікна~$\w$, частки несправних
$\phi$, виробничого масштабування $\kappaScale>0$ від одиниць \RNM{} до
абсолютної пропускної здатності та неспадної $\Cost(\n)$ мінімізувати
$\Cost(\n)$ за умов
$\Prob[\kappaScale\cdot\TPSn(\n,\cc^{\star})\ge\Dpk]\ge1-\epsp$ і
$\pmiss(\n,\w)\le\epss$, де $\cc^{\star}=\arg\max\gfun$. Якщо $\kappaScale$
вільний, брати найменше значення, що задовольняє обмеження пропускної здатності
при $\n=4$.
\end{Problem}

Імовірнісне обмеження еквівалентне, через прогнозний інтервал вище, умові
\begin{equation}
h(\n)\eqdef\log(\kappaScale\,\Tzero\,\gfun(\cc^{\star})/\Dpk)
-\beta\log\n-\zq_{1-\epsp}\hat{V}^{1/2}(\n,\cc^{\star})\ge0.
\label{eq:h}
\end{equation}
На діапазоні модель-базованої екстраполяції ($\n\ge\resNmaxMeasured$) підігнана
коваріація робить $\hat{V}$ неспадною за $\log\n$, тож при $\beta>0$
відображення $h$ там строго спадне; на виміряній сітці $h$ обчислюємо
безпосередньо.

\begin{Proposition}[Вікно допустимості]
\label{prop:window}
Коли $h$ строго спадне і використано наближений поріг детекції
\eqref{eq:nmin}, допустима множина є
$\{\n:\nmin\le\n\le\nmax\}$ з $\nmin$ з~\eqref{eq:nmin} і
$\nmax=\lfloor\sup\{\n:h(\n)\ge0\}\rfloor$. Якщо вона непорожня і $\Cost$
неспадна, то $\nopt=\nmin$.
\end{Proposition}

\begin{Proof}
Детекція виконується на $\{\n\ge\nmin\}$, пропускна здатність --- на
$\{\n\le\nmax\}$, отже інтервал; монотонна вартість вибирає лівий кінець.
\end{Proof}

\begin{Corollary}[Порожня допустима множина]
\label{cor:infeasible}
Якщо $\nmin>\nmax$, жодна кількість валідаторів не задовольняє обидві вимоги за
зазначених обмежень: потрібно змінювати архітектуру, а не лише забезпечення.
\end{Corollary}

Протокольно ефективні розміри --- $\n=3f+1$; інші кількості присутні на сітці
вимірювань, але рішення про закупівлю перетинає вікно з $\{3f+1\}$. Вікно ---
перетин двох монотонних обмежень
\cite{charnes1959,nemirovski2007}. У запропонованій методиці воно перетворює
калібровану модель на відтворюваний робочий процес визначення розміру для
дослідженого класу розгортань.
